\documentclass[10pt,a4paper]{amsart}
\usepackage[margin=19mm]{geometry}
\usepackage{amsmath,amssymb,mathtools}
\usepackage[scr=rsfs,cal=euler]{mathalfa}
\usepackage{xfrac}
\usepackage[unicode,hidelinks,bookmarks=false]{hyperref}
\newcommand{\ie}{i.e.\ }
\newcommand{\cf}{cf.\ }
\newcommand{\resp}{respectively\ }
\newcommand{\Subsection}{\subsection}
\providecommand{\nobreakdash}{\nobreak\hbox{-}}
\setlength{\emergencystretch}{2em}
\multlinegap0pt
\usepackage{bm}
\usepackage{bbm}
\usepackage{dsfont}
\usepackage{xspace}
\usepackage{cancel}
\usepackage{mathtools}
\usepackage{amsthm}
\makeatletter
\@ifundefined{thm}{\newtheorem{thm}{Thm}}{}
\@ifundefined{lem}{\newtheorem{lem}[thm]{Lemma}}{}
\@ifundefined{obs}{\newtheorem{obs}[thm]{Observation}}{}
\makeatother
\newcommand{\ff}{\mathcal{F}}
\newcommand{\m}{\mathcal}
\renewcommand{\ge}{\geqslant}
\title[Delta-system method: a survey]{Delta-system method: a survey}
\author{Andrey Kupavskii}
\date{Source: arXiv:2508.20132v1 (CC BY 4.0)}
\begin{document}
\maketitle

\noindent\emph{Source and licence.} Delta-system method: a survey. Version arXiv:2508.20132v1 (CC BY 4.0), distributed under the Creative Commons Attribution 4.0 International licence, as recorded in the arXiv licence metadata for this exact version and shown on the abstract page https://arxiv.org/abs/2508.20132v1. Full rights notice: licenses/arxiv-2508.20132-CC-BY-4.0.txt.\par\smallskip

\noindent\textit{Editorial scope.} Counted unit: the Lem whose source label is \texttt{lemfur1} in the article (source0.tex lines 482--482, source label \texttt{lemfur1}), delivered together with the article's own complete original proof of it (source0.tex lines 483--498, one \texttt{proof} environment of 16 lines). No mathematics is written, repaired or re-derived: statement and proof are the article's own text line for line. Required context retained uncounted: obsdelta2 (lines 285--286). Every definition, display and label of those lines is retained unchanged. Omitted from this package and not redistributed: 1--284, 287--481, 499--1650. Nothing inside a retained excerpt is omitted or altered: every retained body is delivered exactly as the article's lines are written, so no omission receipt of any kind accompanies this package. Citations: the retained excerpts carry no citation command at all, so no bibliography entry of the article is transcribed and no reference is redistributed. Reference adaptation: none was needed and none was made; every reference inside the delivered unit resolves inside it, so no unresolved reference or citation is produced. Printed slips kept literal: no printed word, symbol, hypothesis or number of the counted statement or of its proof was altered. Layout and class: the article's own class is \texttt{svmult} with packages including type1cm, makeidx, graphicx, multicol, footmisc, newtxtext, newtxmath, hyperref, cprotect; the wrapper is the lane's pinned \texttt{amsart} editorial wrapper at 10pt on A4, which restates the theorem-like environments the retained text needs and adds the article's own macro definitions that the retained text needs (\texttt{\textbackslash ff}, \texttt{\textbackslash ge}, \texttt{\textbackslash m}), with extra wrapper packages (bm, bbm, dsfont, xspace, cancel, mathtools). The delivered numbering is the unit's own. Assets: no retained span contains a figure environment, a graphics-inclusion command, a PDF-page inclusion, a table or a TikZ picture; no image file is copied into this package, the package declares no asset dependency and no image file is redistributed with it. Licence: version arXiv:2508.20132v1 of this article is distributed under the Creative Commons Attribution 4.0 International licence, as recorded in the arXiv licence metadata for this exact version and shown on the abstract page https://arxiv.org/abs/2508.20132v1. A full rights notice is delivered at licenses/arxiv-2508.20132-CC-BY-4.0.txt. Characters: every retained excerpt is pure ASCII, so the delivered engine, XeLaTeX with its default Latin Modern text font, reports no missing character.
\begin{obs}\label{obsdelta2} Let $k$ be an integer. Let $A$ be set of size $k$ and $\m F$ is a $\Delta(k+1)$-system with kernel $D$. Then $|A\cap D| = |A\cap F|$ for some $F\in \m F$.
\end{obs}

    \begin{lem}\label{lemfur1} Assume that $\ff$ does not satisfy the conditions of the theorem. Then we can find $\ff''\subset \ff$ such that $$|\ff''|\ge \frac{|\ff|}{1+s(k-1)|\m B(\ff)|}\ge \frac{|\ff|}{1+s(k-1)2^{k}}$$ and such that either $\ff''$ satisfies the conditions of the theorem or such that there is a set $I\in \m M(\ff)\setminus \m M(\ff'')$.\end{lem}

    \begin{proof}
      Let us define a partition
      $$\ff = \ff_0\cup \bigcup_{I\subset [k]}\ff_I$$
      as follows. At each iteration we go over all sets $A$ of $\ff$ and check if for every $X\subset A$ such that $\pi(X)\in \m M(\ff)$ the set $X$ is a kernel of a $\Delta(s)$-system. If it is not, then we delete this set $A$ from $\ff$ and include it into $\ff_{\pi(X)}$. At some point the procedure stops and we are left with a family $\ff_0$ in which each such $X\subset A$ is a center of an $s$-sunflower for every $A$.

      Next, we note two things. First, we claim that $\ff_0$ is a $(k,s)$-homogeneous structure. We see that $\ff_0$ satisfies conditions (2) and (4) of the theorem by definition. Let us also check that (3) holds for $\ff_0$. Suppose that, for some $A$ and $B$, there is $I\in \ \pi(\ff_0|_A)\setminus \pi(\ff_0|_B)$. Then take the set $X\subset B$ such that $\pi(X)=I$ and note that $X$ is not a kernel of an $\Delta(s)$-system $F_1,\ldots, F_s$: otherwise, by Observation~\ref{obsdelta2}, there would have been $i$ such that $F_i\cap B = X$ since $s>k$. Since $X$ is not a kernel of a $\Delta(s)$-system, such $B$ should have then been `filtered' into $\ff_{\pi(X)}$ during the initial procedure.\footnote{Note here the delicate detail of the initial procedure of `filtering' sets. We check the `kernel condition' for all $X\subset A$ such that $\pi(X)\in \m M(\ff)$, and not only, say, $\pi(X)\in \pi(\ff|_A)$.}  Therefore, $\pi(\ff_0|_A)= \pi(\ff_0|_B)$ for all $A,B$ and $\ff_0$ is as desired.


       Second, we claim that, for each $I\subset [k]$, $I\notin \m B(\ff_{I})$. Indeed, take any $X$ such that $\pi(X) = I$ and look at the first set $F$ containing $X$ that was included in $\ff_I$. At that point, we know that $X$ is not a center of a $\Delta(s)$-system in the remaining part of $\ff$. Only sets from the remainder of $\ff$ may end up in $\ff_I$, and thus $X$ will not become such a center.

       Now, for each $I\subset [k]$ and $X$ such that $\pi(X) = I$ and $\ff_I[X]$ is non-empty, we take a maximal $\Delta$-system $F_1,\ldots, F_\ell\in \ff_I$ with kernel $X$, where $\ell<s$. Any other set in $\ff_I[X]$ (the sets in $\ff_I$ containing X) must intersect one of $F_i\setminus X$, so we may take the element $v_X\in \cup F_i\setminus X$ that is contained in at least a $1/\ell k$-proportion of sets from $\ff_I[X]$. Recall that $\ff$ is $k$-partite, and thus $\ff[X_1]$, $\ff[X_2]$ are disjoint for distinct $X_1,X_2$ such that $\pi(X_i) = I$. Now put $$\ff'_I:=\bigsqcup_{X: \pi(X)=I} \ff_I[X\cup \{v_X\}].$$ We note that, first, $|\ff'_I|\ge \frac 1{(s-1)k}|\ff_I|$ and, second, $I\notin \m M(\ff'_I)$.



      Finally, we take as $\ff''$ the largest family out of  $\ff_0$ and $\{\ff_{I}'\}$. Since $\ff_0$ and $\{\ff_{I}'\}$ decompose $\ff$, we can see that one of the families in question has size at least $\frac{|\ff|}{1+s(k-1)|\m B(\ff)|}$.
    \end{proof}

\end{document}
